% SPDX-License-Identifier: 0BSD
\section{Prepared-Model Eager DAG Execution}
\label{sec:eager-execution}

This section compares the compiled and eager realizations of the shared-current
DAG; recurrence execution is described by the general recursion in
Sec.~\ref{sec:zgg-example}.  Compiled mode combines many current relations into
process-specific stage evaluators.  Eager mode prepares model-local kernels
once and represents each process by an invocation schedule.  They use the same
physical graph, normalization, colour contraction, and resolved observables.

\begin{center}
\small
\begin{tabular}{L{0.18\textwidth}L{0.35\textwidth}L{0.35\textwidth}}
\toprule
 & \textbf{Compiled mode} & \textbf{Eager mode} \\
\midrule
Numerical unit
  & Multi-output expression specialized to one process stage
  & Reusable model kernel applied by a process schedule \\
Process preparation
  & Builds and compiles process-specific stage evaluators
  & Maps the proven DAG to reusable model kernels and a process schedule \\
Runtime work
  & Calls a small number of wide stage evaluators
  & Applies reusable kernels and combines the resulting currents \\
Expected trade-off
  & More preparation for greater process-wide fusion
  & Faster process preparation with additional scheduling work \\
\bottomrule
\end{tabular}
\end{center}

\subsection{Prepared model and process schedule}

A prepared model separates model-level algebra from process-level graph
construction.  Schematically,
\begin{equation}
  \mathcal{B}_{M,b,o,t}
  = \left(
      \mathcal{I}_{M},\;
      \{K_\kappa^{(b,o,t)}\}_{\kappa\in\mathcal{K}_M},\;
      \mathcal{P}_{M}
    \right),
  \label{eq:eager-prepared-model}
\end{equation}
where \(\mathcal I_M\) defines the model, \(K_\kappa\) are numerical kernels for
its canonical interaction classes, and \(\mathcal P_M\) records the numerical
and target conventions.  The labels \(b,o,t\) identify the backend,
optimization, and target architecture.

For a process \(P\), the eager schedule records the graph dependencies, the
kernel associated with each interaction, the current transformations, and the
final colour and helicity reductions. It also retains the dependencies needed
for resolved helicity and colour-flow selections.

Because the numerical kernels are already part of \(\mathcal B\), eager
process generation does not compile a new evaluator for every process.  Model
preparation occurs before the process-generation timer and is not included in
the process-generation values reported here.

\subsection{Stage execution}

Let \(V_{c\alpha p}\) denote component \(\alpha\) of current \(c\) at
phase-space point \(p\).  For one stage, unpropagated components are
\begin{equation}
  U_{c\alpha p}^{(s)} =
  \sum_{r\in\mathcal{R}_s(c,\alpha)}
    w_r\,
    \left[
      K_{\kappa(r)}
      \left(
        G_r(V_{\cdot\cdot p}),\;q_r(p),\;\theta
      \right)
    \right]_{\alpha(r)} ,
  \label{eq:eager-invocation-sum}
\end{equation}
where \(G_r\) gathers the required current components and momenta, \(w_r\) is
an exact factor, and \(\theta\) denotes model parameters.  Contributions are
accumulated before finalization:
\begin{equation}
  V_{c\beta p}^{(s)} =
  \left[F_{\phi(c)}\left(U_{c\cdot p}^{(s)},Q_c(p),\theta\right)\right]_\beta .
  \label{eq:eager-finalization}
\end{equation}
The map \(F_{\phi(c)}\) is an identity, propagator, or another model-defined
current transformation.  Applying it once per current preserves the
normalization convention of compiled mode.

After the last stage, closure relations produce amplitude components,
\begin{equation}
  A_{h\gamma p}
  = \sum_{r\in\mathcal{C}(h,\gamma)}
      w_r^{\mathrm{cl}}
      K_{\kappa(r)}^{\mathrm{cl}}
      \left(G_r^{\mathrm{cl}}(V_{\cdot\cdot p}),q_r(p),\theta\right),
  \label{eq:eager-closure}
\end{equation}
and the standard colour/helicity reduction gives the requested observable.  In
a contracted colour basis,
\begin{equation}
  \mathcal{M}_p =
  \mathcal{N}\sum_h\sum_{\gamma,\delta}
  A_{h\gamma p}\,C_{\gamma\delta}\,A^*_{h\delta p},
  \label{eq:eager-reduction}
\end{equation}
with the same colour matrix, symmetry factors, initial-state average, and
coupling normalization as the other execution modes.

\subsection{Batched evaluation and selectors}

Compatible kernel applications are batched across phase-space points so that
the point dimension can be vectorized. Large input batches may be partitioned,
but this changes neither the graph nor the numerical result.

A helicity or flow selection retains the complete dependency closure of the
requested amplitudes. A complete leading-colour process therefore supports a
selected flow with summed helicities, a selected helicity with all flows, or
any retained subset without changing its physics content. The two corresponding
graph layouts are described in Sec.~\ref{sec:lc-flow-layouts}.

\subsection{Precision, portability, and interpretation}

A standalone eager process contains the numerical kernels and schedule needed
for evaluation.  Portable JIT state is compiled for the receiving processor at
load time; target-native ASM and C++ payloads remain architecture-specific.
The exact symbolic representation is retained for higher-precision evaluation.

Eager generation time starts from a prepared model and ends with a validated,
loadable process; prepared-model creation is outside this timing boundary. At
runtime, eager mode avoids process-specific backend compilation but performs
explicit scheduling and reusable-kernel dispatch.
Compiled mode can be faster when process-wide fusion dominates, whereas eager
mode can reduce preparation cost when many processes share one model.  The
tables quantify this trade-off under identical physics and validation
conditions rather than assuming a universal ranking.
